% EP001: Physical System Description -- Unified Learning Module
% Source: LT-7 Paper + E003 Slides + Thesis
%==============================================================================

\documentclass[11pt,a4paper]{article}

% --- Core packages -----------------------------------------------------------
\usepackage[utf8]{inputenc}
\usepackage[T1]{fontenc}
\usepackage{amsmath,amssymb,amsthm,amsfonts}
\usepackage{geometry}
\usepackage{graphicx}
\usepackage{booktabs}
\usepackage{array}
\usepackage{fancyhdr}
% hyperref loaded LAST to avoid numbering conflicts (no titlesec used)
\usepackage[hypertexnames=false,
            colorlinks=true,
            linkcolor=blue,
            filecolor=magenta,
            urlcolor=blue,
            pdftitle={EP001 Physical System Description},
            pdfauthor={Double-Inverted Pendulum Learning Module}]{hyperref}

% --- Page geometry -----------------------------------------------------------
\geometry{left=2.5cm,right=2.5cm,top=2.5cm,bottom=2.5cm,headheight=14pt}

% --- Headers / footers -------------------------------------------------------
\pagestyle{fancy}
\fancyhf{}
\fancyhead[L]{EP001: Physical System Description}
\fancyhead[R]{Learning Module}
\fancyfoot[C]{\thepage}
\renewcommand{\headrulewidth}{0.4pt}

% --- Theorem environments (no optional title arguments) ----------------------
\newtheorem{theorem}{Theorem}[section]
\newtheorem{lemma}[theorem]{Lemma}
\newtheorem{definition}{Definition}[section]
\newtheorem{remark}{Remark}[section]

% --- Document metadata -------------------------------------------------------
\title{\textbf{EP001: Physical System Description}\\[0.5em]
       \large\textit{Double-Inverted Pendulum Control Theory}\\[0.5em]
       \normalsize Episode ID: EP001 \quad Difficulty: Intermediate \quad Duration: 2 h}
\author{Double-Inverted Pendulum Learning Module}
\date{\today}

%==============================================================================
\begin{document}
%==============================================================================

\maketitle
\tableofcontents
\newpage

%------------------------------------------------------------------------------
\section{Introduction}
%------------------------------------------------------------------------------

\subsection{Overview}

The \textbf{double-inverted pendulum (DIP)} is one of the canonical benchmarks
in nonlinear control theory.  A single horizontal force drives a cart along a
track; two rigid links are hinged in series at the cart and must be balanced
upright using that single actuator.

\begin{itemize}
    \item The upright equilibrium is open-loop unstable.
    \item The system is \emph{underactuated}: two rotational degrees of freedom,
          one control input.
    \item Strong nonlinear coupling exists between the two pendulum links.
    \item The system is a standard testbed for sliding mode control (SMC),
          model predictive control (MPC), and reinforcement learning.
\end{itemize}

\subsection{System Characteristics}

\begin{definition}
A mechanical system is \emph{underactuated} if the number of independent
control inputs is strictly less than the number of degrees of freedom.
\end{definition}

The DIP satisfies this definition with three degrees of freedom (cart position,
link-1 angle, link-2 angle) and one control input (cart force).  Key dynamical
properties are listed below.

\begin{enumerate}
    \item \textbf{Nonlinear coupling}: the inertia matrix depends on joint
          angles, creating state-dependent effective mass.
    \item \textbf{Unstable equilibrium}: small perturbations cause the
          pendulums to fall without active feedback.
    \item \textbf{Non-minimum phase tendencies}: zero dynamics near upright may
          exhibit unstable behavior depending on output choice.
    \item \textbf{Disturbance sensitivity}: external impulses propagate rapidly
          through the kinematic chain.
\end{enumerate}

%------------------------------------------------------------------------------
\section{Physical System Description}
%------------------------------------------------------------------------------

\subsection{System Configuration}

The DIP hardware consists of the following components.

\begin{itemize}
    \item \textbf{Cart} (mass $m_0$): moves along a horizontal track;
          travel limited to $\pm 1\,\mathrm{m}$ in simulation.
    \item \textbf{Link 1} (mass $m_1$, length $l_1$): rigid rod, hinged at the
          cart top; angle $\theta_1$ measured from the upright vertical.
    \item \textbf{Link 2} (mass $m_2$, length $l_2$): rigid rod, hinged at the
          tip of Link 1; angle $\theta_2$ measured from the upright vertical.
    \item \textbf{Actuator}: motor applies horizontal force $u$ to the cart.
    \item \textbf{Sensing}: optical encoders measure $x$, $\theta_1$, $\theta_2$;
          velocities are estimated by finite differences.
\end{itemize}

\begin{remark}
Angles are measured from the upright position (counterclockwise positive),
consistent with the Euler-Lagrange derivation in
Section~\ref{sec:model}.
\end{remark}

\begin{center}
\begin{verbatim}
      [m2, l2]   Link 2
          |  theta_2
     [m1, l1]   Link 1
          |  theta_1
  ========================  Track
       [ m0 ]   Cart
        <--- u (control force)
\end{verbatim}
\end{center}

\subsection{Physical Constraints}

Operating bounds that any controller must respect:

\begin{enumerate}
    \item Cart position: $x \in [-1,\, 1]\,\mathrm{m}$
    \item Control force magnitude: $|u| \leq u_{\max} = 50\,\mathrm{N}$
    \item Link angles: unconstrained in simulation; hardware stops near
          $\theta_i = \pm\pi$
\end{enumerate}

The mass ordering $m_0 > m_1 > m_2$ ensures the heavier cart provides the
necessary reaction force for upright balance.  The length ordering $l_1 > l_2$
gives the base link greater control authority.

%------------------------------------------------------------------------------
\section{Mathematical Model}
\label{sec:model}
%------------------------------------------------------------------------------

\subsection{State Vector Definition}

\begin{definition}
The system state vector is
\begin{equation}
    \mathbf{x} =
    \bigl[x,\;\theta_1,\;\theta_2,\;\dot{x},\;\dot{\theta}_1,\;
    \dot{\theta}_2\bigr]^{T} \in \mathbb{R}^{6},
\end{equation}
where $x$ is the cart position (m), $\theta_i$ are link angles (rad) from the
upright, and dots denote time derivatives.
\end{definition}

The generalized coordinate vector is $\mathbf{q} = [x,\,\theta_1,\,\theta_2]^T$.

\subsection{Equations of Motion}

The equations of motion are derived via the \textbf{Euler-Lagrange method}.
Let $L = T - V$ denote the Lagrangian.  Applying
\begin{equation}
    \frac{d}{dt}\!\left(\frac{\partial L}{\partial \dot{q}_i}\right)
    - \frac{\partial L}{\partial q_i} = Q_i, \quad i = 1, 2, 3,
\end{equation}
where the generalized forces are $Q_1 = u$, $Q_2 = Q_3 = 0$, yields the
standard manipulator form
\begin{equation}
\label{eq:eom}
    \mathbf{M}(\mathbf{q})\,\ddot{\mathbf{q}}
    + \mathbf{C}(\mathbf{q},\dot{\mathbf{q}})\,\dot{\mathbf{q}}
    + \mathbf{G}(\mathbf{q})
    + \mathbf{F}\,\dot{\mathbf{q}}
    = \mathbf{B}\,u + \mathbf{d}(t),
\end{equation}
where $\mathbf{F} = \mathrm{diag}(b_0,b_1,b_2)$ is the viscous friction
matrix, $\mathbf{B} = [1,0,0]^T$, and $\mathbf{d}(t)$ represents bounded
external disturbances.

\subsection{Inertia Matrix}

The inertia matrix $\mathbf{M}(\mathbf{q}) \in \mathbb{R}^{3\times3}$ is
symmetric and configuration-dependent:
\begin{equation}
    \mathbf{M}(\mathbf{q}) =
    \begin{bmatrix}
        M_{11} & M_{12} & M_{13} \\
        M_{12} & M_{22} & M_{23} \\
        M_{13} & M_{23} & M_{33}
    \end{bmatrix},
\end{equation}
with elements derived from the kinetic energy:
\begin{align}
    M_{11} &= m_0 + m_1 + m_2, \\
    M_{12} &= (m_1 r_1 + m_2 l_1)\cos\theta_1 + m_2 r_2\cos\theta_2, \\
    M_{13} &= m_2 r_2\cos\theta_2, \\
    M_{22} &= m_1 r_1^2 + m_2 l_1^2 + I_1, \\
    M_{23} &= m_2 l_1 r_2\cos(\theta_1 - \theta_2) + I_2, \\
    M_{33} &= m_2 r_2^2 + I_2,
\end{align}
where $r_i$ is the distance from joint $i$ to the center of mass of link $i$.

\begin{theorem}
The inertia matrix $\mathbf{M}(\mathbf{q})$ is positive definite for all
physically admissible configurations, i.e., $\mathbf{M}(\mathbf{q}) \succ 0$.
\end{theorem}

\subsection{Coriolis and Centrifugal Matrix}

$\mathbf{C}(\mathbf{q},\dot{\mathbf{q}})$ is computed from Christoffel symbols:
\begin{equation}
    C_{ij} = \sum_{k=1}^{3}\Gamma_{ijk}\,\dot{q}_k, \qquad
    \Gamma_{ijk} = \frac{1}{2}\!\left(
        \frac{\partial M_{ij}}{\partial q_k}
      + \frac{\partial M_{ik}}{\partial q_j}
      - \frac{\partial M_{jk}}{\partial q_i}\right).
\end{equation}
This construction ensures $\dot{\mathbf{M}} - 2\mathbf{C}$ is skew-symmetric,
a property used in passivity-based stability proofs.

\subsection{Gravity and Friction Terms}

The gravity vector follows from potential energy
$V = m_1 g r_1\cos\theta_1 + m_2 g(l_1\cos\theta_1 + r_2\cos\theta_2)$:
\begin{equation}
    G_1 = 0, \qquad
    G_2 = -(m_1 r_1 + m_2 l_1)g\sin\theta_1, \qquad
    G_3 = -m_2 r_2 g\sin\theta_2.
\end{equation}

%------------------------------------------------------------------------------
\section{Nonlinearity Analysis}
%------------------------------------------------------------------------------

\subsection{Configuration-Dependent Inertia}

The off-diagonal elements $M_{12}$, $M_{13}$, $M_{23}$ depend on trigonometric
functions of the joint angles:

\begin{enumerate}
    \item $M_{12}$ varies by up to 40\% as $\theta_1$ sweeps from $0$ to
          $\pi/4$\,rad (nominal parameters).
    \item $M_{23}$ varies by up to 35\% as $(\theta_1 - \theta_2)$ changes,
          representing coupling between the two links.
    \item Near collinear configurations the condition number of
          $\mathbf{M}(\mathbf{q})$ increases, potentially causing ill-conditioning
          in the numerical inversion required to compute $\ddot{\mathbf{q}}$.
\end{enumerate}

Gains tuned at the upright equilibrium $\mathbf{q} = \mathbf{0}$ may therefore
be suboptimal at large angles, motivating gain scheduling or robust designs.

\subsection{Trigonometric Nonlinearity}

\begin{theorem}
The gravity vector $\mathbf{G}(\mathbf{q})$ is globally bounded:
$\|\mathbf{G}(\mathbf{q})\| \leq G_{\max}$ for all $\mathbf{q}$, where
$G_{\max}$ depends only on system masses, link lengths, and $g$.
\end{theorem}

Dominant nonlinear terms are $\sin\theta_i$ in $\mathbf{G}$ and
$\cos(\theta_1 - \theta_2)$ in $M_{23}$.  For small angles
($|\theta_i| \leq 0.3\,\mathrm{rad}$) the approximations
$\sin\theta_i \approx \theta_i$ and $\cos\theta_i \approx 1$ reduce the
system to a linear time-invariant plant suitable for LQR or pole-placement
design.

%------------------------------------------------------------------------------
\section{System Parameters}
%------------------------------------------------------------------------------

\subsection{Nominal Parameter Values}

Table~\ref{tab:params} lists the nominal parameters used throughout this
learning module and in the accompanying simulation code.

\begin{table}[htbp]
\centering
\caption{Nominal DIP System Parameters}
\label{tab:params}
\begin{tabular}{@{}llcc@{}}
\toprule
Parameter                  & Symbol      & Value  & Unit              \\
\midrule
Cart mass                  & $m_0$       & 1.50   & kg                \\
Link 1 mass                & $m_1$       & 0.20   & kg                \\
Link 2 mass                & $m_2$       & 0.15   & kg                \\
Link 1 length              & $l_1$       & 0.40   & m                 \\
Link 2 length              & $l_2$       & 0.30   & m                 \\
Link 1 COM distance        & $r_1$       & 0.20   & m                 \\
Link 2 COM distance        & $r_2$       & 0.15   & m                 \\
Link 1 moment of inertia   & $I_1$       & 0.0081 & kg\,m$^2$         \\
Link 2 moment of inertia   & $I_2$       & 0.0034 & kg\,m$^2$         \\
Gravitational acceleration & $g$         & 9.81   & m\,s$^{-2}$       \\
Maximum control force      & $u_{\max}$  & 50     & N                 \\
Sampling period            & $\Delta t$  & 0.01   & s                 \\
\bottomrule
\end{tabular}
\end{table}

The inertia values satisfy the parallel-axis theorem constraint
$I_i \geq m_i r_i^2$ (equality for uniform rods).

\subsection{Parameter Uncertainty}

Real hardware deviates from nominal values due to:

\begin{itemize}
    \item Manufacturing tolerances in link lengths ($\pm 2\,\mathrm{mm}$ typical).
    \item Mass variation with attached payload or wear.
    \item Friction coefficients that change with temperature and lubrication state.
    \item Encoder resolution limiting velocity-estimation accuracy.
\end{itemize}

Robust SMC designs must guarantee stability across the uncertainty polytope
$\mathcal{P} = \{\,p \mid p_{\min} \leq p \leq p_{\max}\,\}$ defined by
worst-case parameter deviations.

%------------------------------------------------------------------------------
\section{Linearized Model}
%------------------------------------------------------------------------------

\subsection{Equilibrium Points}

\begin{definition}
A state $\mathbf{x}_e$ is an \emph{equilibrium} of~(\ref{eq:eom}) if
$\dot{\mathbf{q}}_e = \mathbf{0}$, $\ddot{\mathbf{q}}_e = \mathbf{0}$,
and $u_e = 0$.
\end{definition}

Two equilibria are of practical interest:

\begin{enumerate}
    \item \textbf{Upright (unstable)}: $\mathbf{x}_e = \mathbf{0}$ --
          the control objective of this module.
    \item \textbf{Downward (stable)}: $\theta_1 = \theta_2 = \pi$ --
          natural rest position; swing-up controllers target this as a
          starting condition before handing off to the balancing controller.
\end{enumerate}

\subsection{Linearized State-Space Representation}

Linearizing~(\ref{eq:eom}) about $\mathbf{x}_e = \mathbf{0}$ via a
first-order Taylor expansion gives
\begin{equation}
    \dot{\mathbf{x}} = \mathbf{A}\,\mathbf{x} + \mathbf{b}\,u,
\end{equation}
where $\mathbf{A} \in \mathbb{R}^{6\times6}$ and $\mathbf{b} \in \mathbb{R}^{6}$
are constant matrices evaluated at $\mathbf{q} = \mathbf{0}$.

Numerical analysis with the nominal parameters of Table~\ref{tab:params} gives:

\begin{itemize}
    \item Two open-loop eigenvalues with positive real parts (unstable pendulum
          modes).
    \item Four eigenvalues with negative real parts (stable cart and damped
          oscillatory modes).
    \item Controllability matrix rank equal to 6, confirming full controllability.
\end{itemize}

%------------------------------------------------------------------------------
\section{Numerical Properties}
%------------------------------------------------------------------------------

\subsection{State Scaling}

Raw state units span several orders of magnitude: position in metres, angles in
radians, velocities up to $\sim\!10\,\mathrm{rad/s}$ during aggressive
manoeuvres.  Normalized coordinates
\begin{equation}
    \tilde{x} = \frac{x}{x_{\max}}, \quad
    \tilde{\theta}_i = \frac{\theta_i}{\theta_{\max}}, \quad
    \tilde{\dot{x}} = \frac{\dot{x}}{\dot{x}_{\max}},
\end{equation}
improve the conditioning of PSO cost functions and observer design.

\subsection{Time Scaling}

\begin{enumerate}
    \item The highest linearized natural frequency sets the Nyquist requirement;
          $\Delta t = 0.01\,\mathrm{s}$ (100\,Hz) provides adequate margin for
          nominal parameters.
    \item Fixed-step Runge-Kutta 4 integration with this timestep yields
          energy-conservation error below 0.1\% over 10\,s with zero friction
          and zero control input.
    \item Variable-step integrators should use absolute and relative tolerances
          of $10^{-6}$ to resolve fast pendulum modes correctly.
\end{enumerate}

%------------------------------------------------------------------------------
\section{Summary and Key Takeaways}
%------------------------------------------------------------------------------

\subsection{System Overview}

The DIP is a 6-state, 1-input nonlinear plant.  Its equations of motion have
the standard manipulator form
$\mathbf{M}\ddot{\mathbf{q}} + \mathbf{C}\dot{\mathbf{q}}
+ \mathbf{G} = \mathbf{B}u$, which is directly exploited by energy-based and
sliding-mode controllers.  The system is controllable, unstable in open loop,
and well-suited as a benchmark for advanced control algorithms.

\subsection{Key Challenges}

\begin{enumerate}
    \item \textbf{Nonlinearity}: inertia and gravity terms are angle-dependent;
          linear designs lose validity beyond $\pm 0.3\,\mathrm{rad}$.
    \item \textbf{Coupling}: cart motion excites both links; second-link motion
          back-drives the first through $M_{23}$.
    \item \textbf{Underactuation}: one force must simultaneously regulate three
          coordinates, requiring exploitation of internal dynamics.
    \item \textbf{Parameter uncertainty}: friction and mass vary in hardware;
          robust designs must cover the full uncertainty polytope.
\end{enumerate}

\subsection{Control Implications}

\begin{theorem}
For system~(\ref{eq:eom}) with $\mathbf{M}(\mathbf{q}) \succ 0$ and bounded
disturbance $\|\mathbf{d}(t)\| \leq D$, a sliding-mode controller
$u = -\eta\,\mathrm{sgn}(\sigma)$ with gain $\eta > D / \lambda_{\min}(\mathbf{M})$
drives the sliding variable $\sigma$ to zero in finite time.
\end{theorem}

\begin{itemize}
    \item Linear controllers (LQR, PID) are effective near the upright
          equilibrium but require gain scheduling for large-angle operation.
    \item Sliding-mode and adaptive controllers provide robustness to parameter
          uncertainty and bounded disturbances.
    \item PSO-tuned gains consistently improve ISE, IAE, and chattering index
          over manually tuned baselines.
    \item Hardware implementations require constraint-aware tuning to respect
          force limits and cart travel bounds.
\end{itemize}

%==============================================================================
\end{document}
%==============================================================================
